% ============================================================
% 正文：Discussion + Conclusion + Acknowledgement
% 作者：Kehui Chen, Jacky Keung, Xiaoxue Ma
% 主题：MCP 故障定位（fault localization）根因实证研究
% ============================================================

\section{Discussion}\label{sec:discussion}

\subsection{Research Implications for the Research Community}

First, our four-layer taxonomy and the open \textbf{real-bug dataset} complement existing MCP
resources and provide reproducible fault cases for reliability evaluation. Our experiments
reveal a \emph{three-channel exposure model}, namely pushed-to-Server, silent failure,
and caller-exception, demonstrating that vanilla MCP SDKs lack sufficient observability
support, which partly explains the high fault-localisation burden reported by industrial
practitioners. Future research can build upon this dataset to develop automated fault
diagnosis and root-cause inference approaches for MCP systems. Our repair experiment further
shows that the difficulty is concentrated in \emph{finding} the fault: an LLM given the full
source repairs 69.6\% of the auto-verifiable cases regardless of our diagnostic clue, so research
effort should target \emph{localisation} and root-cause inference rather than code-fixing
heuristics.

Second, MCP research should move beyond benchmarks and security attacks to study
\emph{ordinary runtime observability}. We show that MCP faults surface through qualitatively
different channels \emph{simultaneously}, a defining property that existing taxonomies do not
capture. Future fault taxonomies and diagnosis tools should be \emph{layer-aware} and account
for this multi-channel nature, rather than assuming errors are uniformly reported at a single
layer.

Third, the prevalence of silent failure (11/49 = 22.4\% of the reproduced real bugs)
motivates research into \emph{deterministic} silent-failure detection, which is more critical
than handling intermittent cases, since the SDK always hides them.

\subsection{Practical Guidance to Industry}

\paragraph{SDK engineering.} Developers should enrich error metadata and, where the true
origin is known, tag it directly so that errors are not uniformly surfaced as Server-side
exceptions. Consistent with our three-channel model, robust MCP error handling must address all
three channels: (i) prevent pushed-to-Server misattribution by surfacing the channel (or the
true origin layer when determinable); (ii) introduce explicit warning flags for potential
silent-failure conditions so that monitoring pipelines can capture logical misbehaviours that
would otherwise remain hidden; and (iii) ensure caller-side Python exceptions (e.g.\
cancellation, session-lifecycle violations) are surfaced through the MCP error protocol with
stable, structured metadata rather than as ad-hoc runtime tracebacks. Our result shows that
even \emph{without} exposing ground truth, packaging the exposure channel into directional
guidance (the diagnostic-clue tool) measurably improves localisation, so SDKs should emit
channel-aware guidance that points a debugger toward the request rather than only the reported
layer.

\paragraph{Deployment.} Enterprises should formulate unified MCP workflows covering tool
registration, protocol adaptation, permission configuration, and fault troubleshooting
\cite{domanski2023digitalization}, clarifying role boundaries across Host, Client, Server,
and end users. A knowledge-management platform is recommended to aggregate documentation,
FAQs, and troubleshooting cases. Dedicated monitoring systems should be deployed to shorten
fault-recovery time, and standardised invocation logic reduces onboarding cost.

\subsection{Threats to Validity}

\textbf{Interview validity.} Interviewees may lack a full understanding of MCP concepts and
modules. We mitigated this by providing an MCP architecture flowchart during interviews and
allowing ``I don't know'' answers. Sampling bias restricts generalisability: participants are
limited to FinTech and Internet practitioners bilingual in Chinese and English with GitHub
open-source experience.

\textbf{Experimental validity.} All experiments were conducted against a single pinned Python
MCP SDK. Different implementations (e.g., TypeScript SDK) may exhibit divergent error-reporting
behaviours, so our observations may not transfer directly. Our dataset is a complete census of
one library's real reported defects but cannot exhaust all corner cases; and the 49 reproduced
bugs are a representative subset of the 342 active defects (the remainder require
authentication servers or other transports outside the study's scope), so our controlled
observations are drawn from that reproducible subset.

\emph{Data-collection honesty.} Our three-channel model is derived from the raw text returned
by the SDK for each reproduced bug, not from any pre-assigned label. The 49 reproduced bugs are
each bound to a real GitHub issue number, so no fault is fabricated, and we report the remaining
342 active defects as confirmed-but-not-reproduced rather than forcing them into a channel. For
the LLM diagnosis baseline we additionally ran a negative-control ablation
(Section~\ref{sec:results}, \emph{Intervention}) that feeds the agent a bare channel \emph{label}
rather than the guidance; it does \emph{not} improve (0.0\% on non-Server bugs, below the 15.4\%
baseline), confirming that the honest 77.6\%/15.4\% figure is not lowered by our metadata
stripping but reflects genuine difficulty, and that the gain comes from the guidance content
rather than any leakable label.

We do not implement a full-fledged fault-diagnosis system; rather, we expose the underlying
SDK behaviours that cause debugging pain and release a reproducible benchmark. We mitigate
this by comparing experimental findings against practitioner interview feedback, where
observations align with reported industrial pain points. Our diagnosis baseline and the
diagnostic-clue tool are deterministic and reproducible; the clue tool is rule-based (it does
not itself depend on the LLM being diagnosed), and we verified it never reveals the origin
layer, so the measured gain is not an artifact of leaking ground truth. Future work may extend
it to a learned LLM-based diagnosis model.

\section{Conclusion}\label{sec:concl}

This paper presents a sequential-exploratory mixed-methods investigation into runtime fault
localisation in Model Context Protocol (MCP) pipelines. Through semi-structured interviews
with 20 practitioners and a real-bug dataset mined from the entire issue tracker of the
official \verb|modelcontextprotocol/python-sdk| --- a census of 1,343 pure issues reduced to
342 active defects, of which 49 are reproduced at runtime over the production Streamable-HTTP
transport --- we quantify why MCP fault localisation is difficult. We define a four-layer
runtime fault taxonomy and empirically derive a \emph{three-channel exposure model} that
characterises how MCP faults actually surface: \emph{pushed-to-Server} (7/13 non-Server bugs),
\emph{silent failure} (11/49 = 22.4\%), and \emph{caller-exception} (3/13), with a further
17/49 = 34.7\% of Server-origin bugs \emph{genericised}, their real root cause replaced by a
generic message. None of the 13 non-Server bugs is localised at its true origin (10/13 = 76.9\%
generalised mismatch). A surface-text LLM diagnostic agent reading only the SDK's reported
error text achieves 77.6\% overall but only \textbf{15.4\%} on the 13 non-Server bugs (0\% on
both Host- and Client-origin faults, misattributed to the Server layer); a runnable
diagnostic-clue tool that classifies the exposure channel and points the debugger in the right
direction --- without revealing the origin layer --- raises non-Server accuracy to
\textbf{23.1\%} (+7.7 pp), whereas a leaky channel-label control does not help (0\%). A
complementary repair experiment confirms that the difficulty lies in localisation rather than
fixing: given the full source an LLM repairs 69.6\% of the auto-verifiable cases with or
without the clue. Our study converts a subjective industrial pain point into objective,
reproducible technical evidence and a concrete observability fix, and releases the real-bug
dataset for future research on MCP fault diagnosis and reliability.

In the future, we plan to extend the dataset with additional MCP SDK implementations (e.g.\
TypeScript), broader production failure modes, and to develop a learned LLM-based
fault-diagnosis model evaluated on this benchmark.

\section*{Acknowledgement}

This work is partially supported by the Research Grants Council (RGC) of the Hong Kong
Special Administrative Region, China under the General Research Fund (UGC/FDS16/\allowbreak E27/25),
Hong Kong Metropolitan University Research Grant (RD/2025/\allowbreak 1.24), and research funds from City
University of Hong Kong (6000871, 9229109, 9229029, 9229192). The authors declare no competing
interests.

\FloatBarrier